%% =================================================================
%% Wei, Mei, Zhao, Xiao, Sun, Zhang, Guo.
%% "Nondestructively identifying the mechanical behavior of soft
%%  tissues using surface deformation with an explicit inverse
%%  approach."
%% Appl. Math. Modelling 134 (2024) 126--147.
%%
%% Reconstruction of page 7 (printed p. 132) as study notes.
%% Prose: extracted verbatim from the PDF text layer via pdftotext.
%% Math:  hand-transcribed from the rendered page image —
%%        Eqs. (18)–(26), nine equations on a single page.
%%        All of these are hand-derived gradient expressions that
%%        autodiff (e.g. PyTorch's loss.backward()) would handle
%%        automatically with zero manual derivation.
%% =================================================================

\documentclass[11pt]{article}

\usepackage[margin=0.9in]{geometry}
\usepackage{amsmath, amssymb}
\usepackage{mathrsfs}
\usepackage{bm}
\usepackage{fancyhdr}

\pagestyle{fancy}
\fancyhf{}
\lhead{\small Y.~Mei et al.}
\rhead{\small Applied Mathematical Modelling 134 (2024) 126--147}
\cfoot{\small \thepage}
\renewcommand{\headrulewidth}{0.4pt}

\setcounter{page}{132}
\setcounter{equation}{17}  % next equation is (18)

\begin{document}

\noindent
nodal or elemental shear moduli. If the shear modulus is nodally
defined, we can use the same interpolation function with the
displacement field. In this case, the shear modulus can be expressed
as:
\begin{equation}
\mu \approx \mu^{h} = \sum_{i=1}^{N} N_i\, \mu_i
\end{equation}
where $N$ is the total number of nodes in the finite element model
and $N_i$ is the finite element shape function. $\mu_i$ is the nodal
shear modulus.

In the explicit inverse method, since shear modulus distribution is a
function of geometric parameters and the shear modulus in each
region, the gradient of the objective function with respect to the
optimization variables can be calculated using the chain rule. To
this end, we should seek $\delta\mu$ corresponding to a change in any
optimization variable from $\beta$ to $\beta + \delta\beta$:
\begin{equation}
\delta\mu = \frac{\partial\mu}{\partial\beta}\,\delta\beta
\end{equation}
Thereby, comparing to the implicit inverse method, we need to
evaluate the $\partial\mu / \partial\beta$.

\medskip
\noindent
\textbf{(a) Layered structure}

\medskip
\noindent
The evaluation of the gradient of geometric optimization variables in
a layered structure with $n + 1$ layers should be discussed
considering three cases. For the $m$-th geometric optimization
variable $r^{1}_{m}$ at the first interface $C^{1}$,
$\frac{\partial\mu}{\partial r^{1}_{m}}$ is given by:
\begin{equation}
\frac{\partial\mu}{\partial r^{1}_{m}}
= \left[-\mu^{0} + \mu^{1}\!\left(1 - H\!\left(\phi^{2}(\mathbf{x})\right)\right)\right]
  \frac{\partial H\!\left(\phi^{1}(\mathbf{x})\right)}{\partial\phi^{1}(\mathbf{x})}\,
  \frac{\partial\phi^{1}(\mathbf{x})}{\partial r^{1}_{m}}
\end{equation}

For the $m$-th geometric optimization variable $r^{n}_{m}$ at the
last interface $C^{n}$,
$\frac{\partial\mu}{\partial r^{n}_{m}}$ is given by:
\begin{equation}
\frac{\partial\mu}{\partial r^{n}_{m}}
= \left(-\mu^{n-1}\, H\!\left(\phi^{n-1}(\mathbf{x})\right) + \mu^{n}\right)
  \frac{\partial H\!\left(\phi^{n}(\mathbf{x})\right)}{\partial\phi^{n}(\mathbf{x})}\,
  \frac{\partial\phi^{n}(\mathbf{x})}{\partial r^{n}_{m}}
\end{equation}

For the $m$-th geometric optimization variable at the intermediate
interface $C^{i}\,(2 \leq i \leq (n-1))$,
$\frac{\partial\mu}{\partial r^{i}_{m}}$ is given by:
\begin{equation}
\frac{\partial\mu}{\partial r^{i}_{m}}
= \left[-\mu^{i-1}\, H\!\left(\phi^{i-1}(\mathbf{x})\right)
        + \mu^{i}\!\left(1 - H\!\left(\phi^{i+1}(\mathbf{x})\right)\right)\right]
  \frac{\partial H\!\left(\phi^{i}(\mathbf{x})\right)}{\partial\phi^{i}(\mathbf{x})}\,
  \frac{\partial\phi^{i}(\mathbf{x})}{\partial r^{i}_{m}}
\end{equation}

In Eqs.~(20--22),
$\frac{\partial H(\phi^{j}(\mathbf{x}))}{\partial\phi^{j}(\mathbf{x})}$
$(1 \leq j \leq n)$ can be evaluated by Eq.~(13) as the Heaviside
function is approximated by a continuous and differentiable function.
$\frac{\partial\phi^{j}(\mathbf{x})}{\partial r^{j}_{m}}$ can be
evaluated by the finite difference method.

For the shear modulus of the background $\mu^{0}$,
$\frac{\partial\mu}{\partial\mu^{0}}$ is given by:
\begin{equation}
\frac{\partial\mu}{\partial\mu^{0}} = 1 - H\!\left(\phi^{1}(\mathbf{x})\right)
\end{equation}

For the shear modulus of the last layer $\mu^{n}$, we have
\begin{equation}
\frac{\partial\mu}{\partial\mu^{n}} = H(\phi^{n}(\mathbf{x}))
\end{equation}

For the shear modulus of the layer $\mu^{i}\,(1 \leq i \leq (n-1))$,
we have
\begin{equation}
\frac{\partial\mu}{\partial\mu^{i}}
= H\!\left(\phi^{i}(\mathbf{x})\right)
  \left(1 - H\!\left(\phi^{i+1}(\mathbf{x})\right)\right)
\end{equation}

\medskip
\noindent
\textbf{(b) Domain with inclusion embedded}

\medskip
\noindent
For each $m$-th geometric optimization variables of $j$-th component
$r^{j}_{m}$,
$\frac{\partial\mu}{\partial r^{j}_{m}}$ is given by:
\begin{equation}
\frac{\partial\mu}{\partial r^{j}_{m}}
= \left[\mu^{0}\!\prod_{\substack{i=1 \\ i \neq j}}^{n}
         H\!\left(\phi^{i}(\mathbf{x})\right)
  - \mu^{j}\right]
  \frac{\partial H\!\left(\phi^{j}(\mathbf{x})\right)}{\partial\phi^{j}(\mathbf{x})}\,
  \frac{\partial\phi^{j}(\mathbf{x})}{\partial r^{j}_{m}}
\end{equation}

For the shear modulus of the background $\mu^{0}$, we have:

% [page ends here — continues on page 8]

\end{document}
